
\chapter*{Expos\'e I\\Algebraic Structures. Group Cohomology}
\addcontentsline{toc}{chapter}{Expos\'e I. Algebraic Structures. Group Cohomology}
\markboth{EXPOS\'E I. ALGEBRAIC STRUCTURES. GROUP COHOMOLOGY}%
  {EXPOS\'E I. ALGEBRAIC STRUCTURES. GROUP COHOMOLOGY}

\begin{center}
  \textit{by M. Demazure}
\end{center}

This expos\'e consists of two parts: the first collects a number of general
definitions and establishes notation that will often be used in the sequel;
the second treats group cohomology and culminates in \SGARefLink{sga3:i:theorem:5-3-3}{Theorem~5.3.3}, on the
vanishing of the cohomology of diagonalizable groups.

We choose once and for all a Universe.\footnote{Translator's note: on universes, see SGA~4, Exposé~I, \S0 and the
Appendix. The Gille--Polo re-edition (2011), Exposé~I, gives this reference
and also refers to the discussion in M.~Demazure and P.~Gabriel,
\textit{Groupes algébriques}, tome~I, Masson and North-Holland, 1970,
p.~xxvi.}
All definitions and constructions will be relative to this Universe. We shall
systematically allow ourselves the following abuse of language. To define a
functor \(f:\mathscr C\to\mathscr C'\), we shall be content to define the
object \(f(S)\) of \(\mathscr C'\) for every object \(S\) of \(\mathscr C\),
whenever there is no ambiguity about how to define \(f(h)\) for an arrow \(h\)
of \(\mathscr C\). In practice, we shall say: let
\(f:\mathscr C\to\mathscr C'\) be the functor defined by
\(f(S)=\cdots\).

\SGARefTarget{sga3:i:section:1}
\section*{1. Generalities}
\addcontentsline{toc}{section}{1. Generalities}

\subsection*{\SGAInlineRefTarget{sga3:i:subsection:1-1}1.1.}

Let \(\mathscr C\) be a category. We denote by \(\widehat{\mathscr C}\) the
category
\[
  \operatorname{Hom}\bigl(\mathscr C^{\circ},\mathbf{Set}\bigr)
\]
of contravariant functors from \(\mathscr C\) to the category \(\mathbf{Set}\)
of sets. There is a canonical functor
\[
  \mathbf h:\mathscr C\longrightarrow\widehat{\mathscr C}
\]
which associates with every \(X\in\operatorname{Ob}(\mathscr C)\) the functor
\(\mathbf h_X\) such that
\[
  \mathbf h_X(S)=\operatorname{Hom}(S,X).
\]
For every functor \(\mathbf F\in\operatorname{Ob}(\widehat{\mathscr C})\),
one has (see, for example, EGA~\(0_{\mathrm{III}}\), 8.1) a bijection
\[
  \operatorname{Hom}(\mathbf h_X,\mathbf F)
    \xrightarrow{\ \sim\ } \mathbf F(X).
\]

In particular, for every pair \(X,Y\) of objects of \(\mathscr C\), the
canonical map
\[
  \operatorname{Hom}_{\mathscr C}(X,Y)
    \xrightarrow{\ \sim\ }
  \operatorname{Hom}_{\widehat{\mathscr C}}
    (\mathbf h_X,\mathbf h_Y)
\]
is bijective; the functor \(\mathbf h\) is fully faithful. It
therefore defines an isomorphism of \(\mathscr C\) onto a full subcategory of
\(\widehat{\mathscr C}\), and an equivalence of \(\mathscr C\) with the full
subcategory of \(\widehat{\mathscr C}\) formed by the representable functors
(that is, the functors isomorphic to one of the form \(\mathbf h_X\)). In the
sequel we shall often identify \(X\) with \(\mathbf h_X\). The following
paragraphs are intended to show that this identification can be made safely.\footnote{Translator's note: the Gille--Polo re-edition (2011), Exposé~I, adds here a
variant. In our formulation: if \(\mathscr D\) is a full subcategory of
\(\mathscr C\) containing \(X\) and \(Y\), and \(\mathbf h'_X\),
\(\mathbf h'_Y\) are the restrictions of \(\mathbf h_X\), \(\mathbf h_Y\) to
\(\mathscr D\), the bijection above, applied in \(\widehat{\mathscr D}\),
gives \(\operatorname{Hom}(\mathbf h'_X,\mathbf h'_Y)\simeq\mathbf h'_Y(X)
=\operatorname{Hom}(X,Y)\). Thus a morphism \(X\to Y\) is the same as a
family of maps \(\operatorname{Hom}(T,X)\to\operatorname{Hom}(T,Y)\),
\(T\in\operatorname{Ob}(\mathscr D)\), functorial in \(T\).}

\subsection*{\SGAInlineRefTarget{sga3:i:subsection:1-2}1.2.}

A monomorphism in \(\widehat{\mathscr C}\) is precisely a morphism
\(\mathbf F\to\mathbf G\) such that, for every
\(S\in\operatorname{Ob}(\mathscr C)\), the corresponding map of sets
\(\mathbf F(S)\to\mathbf G(S)\) is injective.\footnote{Translator's note: the Gille--Polo re-edition (2011), Exposé~I, adds a proof and
the corresponding statement for epimorphisms. In our formulation: a morphism
\(u:\mathbf F\to\mathbf G\) injective at every \(S\) is clearly a
monomorphism; conversely, if \(u\) is a monomorphism and
\(x,y\in\mathbf F(S)\) have the same image in \(\mathbf G(S)\), the
morphisms \(\mathbf h_S\to\mathbf F\) defined by \(x\) and \(y\) (1.1) have
the same composite with \(u\), hence are equal, and \(x=y\). Likewise,
\(u\) is an epimorphism if and only if every \(u(S)\) is surjective: the
pushout \(\mathbf G\amalg_{\mathbf F}\mathbf G\) is computed pointwise, and
the two canonical maps from \(\mathbf G(S)\) to
\(\mathbf G(S)\amalg_{\mathbf F(S)}\mathbf G(S)\) agree only if
\(\mathbf F(S)\to\mathbf G(S)\) is surjective.}
We shall say that \(\mathbf F\) is a \emph{subobject} (or a subfunctor) of
\(\mathbf G\) if \(\mathbf F(S)\) is a subset of \(\mathbf G(S)\) for every
\(S\).

In \(\widehat{\mathscr C}\), arbitrary inverse limits exist and are computed
pointwise\footnote{Translator's note: the Gille--Polo re-edition (2011), Exposé~I, adds two
remarks here. Arbitrary direct limits also exist in \(\widehat{\mathscr C}\)
and are computed pointwise, but \(\mathbf h\) does not commute with them in
general: for instance, if \(\varnothing\) is an initial object of
\(\mathscr C\), then \(\mathbf h_\varnothing(\varnothing)\) is not empty, so
\(\mathbf h_\varnothing\) is not the initial object of
\(\widehat{\mathscr C}\). And the kernel of a pair of morphisms
\(u,v:\mathbf F\rightrightarrows\mathbf G\) is the subfunctor
\(S\mapsto\{x\in\mathbf F(S)\mid u(x)=v(x)\}\) of \(\mathbf F\).}:
\[
  \left(\varprojlim_i\mathbf F_i\right)(S)
    =\varprojlim_i\mathbf F_i(S).
\]
In particular, fiber products are defined by
\[
  (\mathbf F\mathbin{\times_{\mathbf G}}\mathbf F')(S)
    =\mathbf F(S)\mathbin{\times_{\mathbf G(S)}}\mathbf F'(S).
\]

As a final object of \(\widehat{\mathscr C}\), we choose the functor
\(\underline{\mathbf e}\) such that
\(\underline{\mathbf e}(S)=\{\varnothing\}\).
Every \(\mathbf F\in\operatorname{Ob}(\widehat{\mathscr C})\) has a unique
morphism to \(\underline{\mathbf e}\), and we set
\[
  \mathbf F\times\mathbf F'
    =\mathbf F\mathbin{\times_{\underline{\mathbf e}}}\mathbf F'.
\]

For every \(\mathbf F\in\operatorname{Ob}(\widehat{\mathscr C})\), set
\[
  \Gamma(\mathbf F)=\operatorname{Hom}
    (\underline{\mathbf e},\mathbf F).
\]
An element of \(\Gamma(\mathbf F)\) is therefore a family
\((\gamma_S)_{S\in\operatorname{Ob}(\mathscr C)}\), with
\(\gamma_S\in\mathbf F(S)\), such that for every arrow
\(f:S'\to S''\) of \(\mathscr C\),
\[
  \mathbf F(f)(\gamma_{S''})=\gamma_{S'}.
\]

The functor \(\mathbf h\) commutes with inverse limits. In particular, for
\(X\times X'\) to exist, where
\(X,X'\in\operatorname{Ob}(\mathscr C)\), respectively for \(\mathscr C\) to
admit a final object \(e\), it is necessary and sufficient that
\(\mathbf h_X\times\mathbf h_{X'}\), respectively
\(\underline{\mathbf e}\), be representable; in that case
\[
  \mathbf h_X\times\mathbf h_{X'}\simeq\mathbf h_{X\times X'},
  \qquad
  \mathbf h_e\simeq\underline{\mathbf e}.
\]
For \(X\in\operatorname{Ob}(\mathscr C)\), set
\(\Gamma(X)=\Gamma(\mathbf h_X)\). If \(\mathscr C\) has a final object
\(e\), then
\[
  \Gamma(X)\simeq\operatorname{Hom}(e,X).
\]

\subsection*{\SGAInlineRefTarget{sga3:i:subsection:1-3}1.3.}

Let \(S\in\operatorname{Ob}(\mathscr C)\). We denote by
\(\mathscr C_{/S}\) the category of objects of \(\mathscr C\) over \(S\):
its objects are the arrows \(f:T\to S\) of \(\mathscr C\), and
\(\operatorname{Hom}(f,f')\) is the subset of
\(\operatorname{Hom}(T,T')\) formed by the arrows \(u\) such that
\(f=f'\circ u\). If \(\mathscr C\) has a final object \(e\), then
\(\mathscr C_{/e}\) is isomorphic to \(\mathscr C\). The category
\(\mathscr C_{/S}\) has a final object, namely the identity arrow
\(S\to S\).

If \(f:T\to S\) is an object of \(\mathscr C_{/S}\), one may form the
category \((\mathscr C_{/S})_{/f}\), which by abuse of language we denote by
\((\mathscr C_{/S})_{/T}\), and there is a canonical isomorphism
\[
  \mathscr C_{/T}\simeq(\mathscr C_{/S})_{/T}.
\]
This construction applies in particular to the categories
\(\widehat{\mathscr C}\), \(\widehat{\mathscr C_{/S}}\), etc.; in particular,
one defines the category \(\widehat{\mathscr C}_{/\mathbf h_S}\).

If \(f:T\to S\) is an object of \(\mathscr C_{/S}\), then \(\Gamma(f)\)
is identified with the set \(\Gamma(T/S)\) of sections of \(T\) over \(S\),
that is, arrows \(S\to T\) that are right inverses to \(f\). Moreover,
\(\mathbf h_f:\mathbf h_T\to\mathbf h_S\) is then an object of
\(\widehat{\mathscr C}_{/\mathbf h_S}\), and
\[
  \Gamma(\mathbf h_f)
    \simeq\Gamma(\mathbf h_T/\mathbf h_S)
    \simeq\Gamma(T/S)
    \simeq\Gamma(f).
\]

\subsection*{\SGAInlineRefTarget{sga3:i:subsection:1-4}1.4.}

We now define an equivalence between the categories
\(\widehat{\mathscr C_{/S}}\) and
\(\widehat{\mathscr C}_{/\mathbf h_S}\). In other words, we shall prove that
to give a functor on the category of objects of \(\mathscr C\) over \(S\) is
``the same thing'' as to give a functor on \(\mathscr C\) equipped with a
morphism to \(\mathbf h_S\).

\medskip
\noindent\textbf{(i) Construction of
\(\alpha_S:\widehat{\mathscr C}_{/\mathbf h_S}
\to\widehat{\mathscr C_{/S}}\).}

First let \(H:\mathbf F\to\mathbf h_S\) be an object of
\(\widehat{\mathscr C}_{/\mathbf h_S}\). We must define a functor
\(\alpha_S(H)\) on \(\mathscr C_{/S}\). For an object
\(f:T\to S\) of \(\mathscr C_{/S}\), define \(\alpha_S(H)(f)\) to be the
inverse image\footnote{Translator's note: the Gille--Polo re-edition (2011), Exposé~I, adds the
representable case. In our formulation: if \(\mathbf F=\mathbf h_X\), then
\(H\) is given by a morphism \(h:X\to S\), the map \(H(T)\) is
\(g\mapsto h\circ g\), and \(\alpha_S(H)(f)\) is the set of the \(g:T\to X\)
with \(h\circ g=f\), that is, \(\operatorname{Hom}_{\mathscr C_{/S}}(f,h)\);
so \(\alpha_S(H)\) is the functor on \(\mathscr C_{/S}\) represented by the
object \(h\) (compare 1.5.1).} of \(f\in\mathbf h_S(T)\) under the map
\[
  H(T):\mathbf F(T)\longrightarrow\mathbf h_S(T).
\]

Next, let \(u:f\to f'\) be an arrow of \(\mathscr C_{/S}\). Then
\(\mathbf F(u):\mathbf F(T')\to\mathbf F(T)\) induces a map from
\(\alpha_S(H)(f')\) to \(\alpha_S(H)(f)\), denoted \(\alpha_S(H)(u)\). It
is immediate that the assignments
\[
  f\longmapsto\alpha_S(H)(f),
  \qquad
  u\longmapsto\alpha_S(H)(u)
\]
define a functor on \(\mathscr C_{/S}\), hence an object of
\(\widehat{\mathscr C_{/S}}\).

Finally, let \(H:\mathbf F\to\mathbf h_S\) and
\(H':\mathbf F'\to\mathbf h_S\) be two objects of
\(\widehat{\mathscr C}_{/\mathbf h_S}\), and let \(U:H\to H'\) be a
morphism in that category:

\begin{center}
  \begin{tikzcd}[column sep=5.2em,row sep=3.2em]
  \mathbf F
    \arrow[rr,"U"]
    \arrow[dr,"H"']
  &&
  \mathbf F'
    \arrow[dl,"H'"]
  \\
  &
  \mathbf h_S
  &
  {}.
\end{tikzcd}

\end{center}

For every \(f:T\to S\), the map
\(U(T):\mathbf F(T)\to\mathbf F'(T)\) induces a map
\[
  \alpha_S(U)(f):\alpha_S(H)(f)
    \longrightarrow\alpha_S(H')(f),
\]
which defines a morphism of functors
\[
  \alpha_S(U):\alpha_S(H)\longrightarrow\alpha_S(H').
\]
It is easy to verify that the assignments
\[
  H\longmapsto\alpha_S(H),
  \qquad
  U\longmapsto\alpha_S(U)
\]
define a functor
\[
  \alpha_S:\widehat{\mathscr C}_{/\mathbf h_S}
    \longrightarrow\widehat{\mathscr C_{/S}}.
\]

\medskip
\SGARefTarget{sga3:i:proposition:1:4:1}
\noindent\textbf{(ii) \SGARefLink{sga3:i:proposition:1:4:1}{Proposition 1.4.1}.}
\textit{The functor
\(\alpha_S:\widehat{\mathscr C}_{/\mathbf h_S}
\to\widehat{\mathscr C_{/S}}\) is an equivalence of categories.}

We indicate only the principle of the construction of a quasi-inverse functor
\[
  \beta_S:\widehat{\mathscr C_{/S}}
    \longrightarrow\widehat{\mathscr C}_{/\mathbf h_S}.
\]
Let \(\mathbf G\) be a functor on \(\mathscr C_{/S}\). For every object
\(T\) of \(\mathscr C\), set
\[
  \beta_S(\mathbf G)(T)
    =\coprod_{f\in\operatorname{Hom}(T,S)}\mathbf G(f),
  \qquad \operatorname{Hom}(T,S)=\mathbf h_S(T).
\]
This defines a functor \(\beta_S(\mathbf G)\) on \(\mathscr C\), equipped
with the evident projection to \(\mathbf h_S\).

\subsection*{\SGAInlineRefTarget{sga3:i:subsection:1-5}1.5.}

The equivalence \(\alpha_S\) commutes with the functors \(\Gamma\). In other
words, if \(H:\mathbf F\to\mathbf h_S\) is an object of
\(\widehat{\mathscr C}_{/\mathbf h_S}\) and \(\alpha_S(H)\) is the
corresponding object of \(\widehat{\mathscr C_{/S}}\), then
\[
  \Gamma(\alpha_S(H))
    \simeq\Gamma(H)
    \simeq\Gamma(\mathbf F/\mathbf h_S).
\]
The equivalence \(\alpha_S\) commutes with the functors \(\mathbf h\):
if \(f:T\to S\) is an object of \(\mathscr C_{/S}\), then
\(\mathbf h_f:\mathbf h_T\to\mathbf h_S\) is an object of
\(\widehat{\mathscr C}_{/\mathbf h_S}\), and its image under \(\alpha_S\)
is precisely \(\mathbf h_{\mathscr C_{/S}}(f)\), where
\[
  \mathbf h_{\mathscr C_{/S}}:
    \mathscr C_{/S}\longrightarrow\widehat{\mathscr C_{/S}}
\]
is the canonical functor.
Consequently:

\medskip
\noindent\SGAInlineRefTarget{sga3:i:proposition:1-5-1}\textbf{Proposition 1.5.1.}
\textit{Let \(H:\mathbf F\to\mathbf h_S\) be an object of
\(\widehat{\mathscr C}_{/\mathbf h_S}\). The functor
\[
  \alpha_S(H):(\mathscr C_{/S})^{\circ}\longrightarrow\mathbf{Set}
\]
is representable if and only if
\(\mathbf F:\mathscr C^{\circ}\to\mathbf{Set}\) is representable. If
\(\mathbf F\simeq\mathbf h_T\), then \(\alpha_S(H)\) is represented by the
object \(T\to S\) of \(\mathscr C_{/S}\).}

The equivalence \(\alpha_S\) is transitive in \(S\). If \(f:T\to S\) is an
object of \(\mathscr C_{/S}\), there is a commutative diagram of
equivalences:

\begin{center}
  \begin{tikzcd}[column sep=4.0em,row sep=3.0em]
  \bigl(\widehat{\mathscr C}_{/\mathbf h_S}\bigr)_{/\mathbf h_T}
    \arrow[r,"\alpha_{S/\mathbf h_T}"]
    \arrow[dr,"\simeq"']
  &
  \bigl(\widehat{\mathscr C_{/S}}\bigr)_{/\mathbf h_T}
    \arrow[r,"\alpha_f"]
  &
  \widehat{(\mathscr C_{/S})_{/T}}
    \arrow[dl,"\simeq"]
  \\
  &
  \widehat{\mathscr C}_{/\mathbf h_T}
    \arrow[r,"\alpha_T"]
  &
  \widehat{\mathscr C_{/T}}
  \qquad ,
\end{tikzcd}

\end{center}

Here \(\alpha_{S/\mathbf h_T}\) denotes, provisionally, the restriction
(see \SGARefLink{sga3:i:subsection:1-6}{1.6}) of the functor \(\alpha_S\) to the objects over
\(\mathbf h_T\).

